\begin{explanationbox}
\textbf{\textcolor{CExplanation}{一句话直觉}}

当 $x$ 非负时，复杂的对数 $\ln(1+x)$ 可以用简单的分式或多项式控制，便于放缩和估计。

\medskip
\textbf{\textcolor{CExplanation}{核心拆解}}
\begin{description}[style=nextline, leftmargin=2.8em, labelsep=0.5em, itemsep=0.45em]
    \item[\textbf{要点一（形式稳健）：}] 双边估计的上下界都是有理分式，结构对称且精度适中，适合直接套用。
    \item[\textbf{要点二（精度可调）：}] 标准下界分母含 $(1+x)$，比常见的 $x-\frac{x^2}{2}$ 更紧；三阶上界则提供更高精度。
\end{description}

\medskip
\textbf{\textcolor{CExplanation}{几何本质（曲线包夹）}}

在 $x\geq0$ 区间上，曲线 $y=\ln(1+x)$ 被 $y=\frac{2x}{x+2}$（下）和 $y=\frac{x(x+2)}{2(x+1)}$（上）夹在中间，且三条曲线在原点相切。

\medskip
\textbf{\textcolor{CExplanation}{代数意义（分式控制）}}

将超越函数 $\ln(1+x)$ 用代数函数（分式或多项式）进行上下逼近，从而将复杂的对数运算转化为简单的代数运算，便于不等式证明。

\medskip
\textbf{\textcolor{CExplanation}{考点价值（放缩证明）}}
\begin{itemize}[leftmargin=*, itemsep=0.5em, topsep=0.4em]
    \item \textbf{考法一（直接套用）：} 在证明题中直接使用这些不等式进行放缩，简化证明过程。
    \item \textbf{考法二（数值估计）：} 用于估计对数式的近似值或比较两个含对数式的大小。
\end{itemize}

\medskip
\textbf{\textcolor{CExplanation}{顿悟点}}

\textbf{遇到 $\ln(1+x)$（$x\geq0$）需要放缩时，根据题目对精度的要求，从这三个不等式中挑选最合适的一个。}

\medskip
\textbf{\textcolor{CExplanation}{使用场景}}
\begin{itemize}[leftmargin=*, itemsep=0.5em, topsep=0.4em]
    \item \textbf{场景一（导数大题）：} 在导数综合题中，用这些不等式替换 $\ln(1+x)$，将超越问题转化为代数问题。
    \item \textbf{场景二（比较大小）：} 比较 $\ln(1+x)$ 与某个多项式或分式的大小时，可利用对应的界直接判断。
\end{itemize}
\end{explanationbox}
